\chapter{张量代数、外代数、对称代数}

回顾第一章中引入的指数记号，我们将频繁使用它（I，§ 7，第 8 号）：

设 \((x_{\lambda})_{\lambda \in L}\) 是环 \(\mathbf{A}\) 的一族两两可交换的元素；对每个具有有限支撑的映射 \(\alpha: \mathbf{L} \to \mathbf{N}\) ，我们将记

\[
x ^ {\alpha} = \prod_ {\lambda \in L} x _ {\lambda} ^ {\alpha (\lambda)}.
\]

若 \(\beta\) 是另一个 \(\mathbf{L}\) 到 \(\mathbf{N}\) 的具有有限支撑的映射，则 \(\alpha + \beta\) 表示映射

\[
\lambda \mapsto \alpha (\lambda) + \beta (\lambda)
\]

把 L 映入 N；带有这个合成法则，由 L 到 N 的具有有限支撑的映射所成的集 \(\mathbf{N}^{(L)}\) 是由 L 导出的自由交换幺半群，而

\[
x ^ {\alpha} x ^ {\beta} = x ^ {\alpha + \beta},
\]

对所有 \(\alpha \in \mathbf{N}^{(L)}\) ，我们记 \(|\alpha| = \sum_{\lambda \in L} \alpha(\lambda) \in \mathbf{N}\) ；于是 \(|\alpha + \beta| = |\alpha| + |\beta|\) ； \(|\alpha|\) 称为“多重指标” \(\alpha\) 的阶。对所有 \(\lambda \in L\) ，用 \(\delta_{\lambda}\) 记 \(\mathbf{N}^{(L)}\) 中使得 \(\delta_{\lambda}(\lambda) = 1\) , \(\delta_{\lambda}(\mu) = 0\) （对 \(\mu \neq \lambda\) ）的元素（克罗内克指标）；诸 \(\delta_{\lambda}\) （对 \(\lambda \in L\) ）是 \(\mathbf{N}^{(L)}\) 中仅有的阶为 1 的元素。给 \(\mathbf{N}^{(L)}\) 赋予由 \(\mathbf{N}^{\mathrm{L}}\) 上的乘积序诱导的序，使得关系 \(\alpha \leqslant \beta\) 等价于“ \(\alpha(\lambda) \leqslant \beta(\lambda)\) （对所有 \(\lambda \in L\) ）”；这时多重指标 \(\lambda \mapsto \beta(\lambda) - \alpha(\lambda)\) 记作 \(\beta - \alpha\) ，因而它是使 \(\alpha + (\beta - \alpha) = \beta\) 的唯一多重指标。对所有 \(\alpha \in \mathbf{N}^{(L)}\) ，只有有限多个多重指标 \(\beta \leqslant \alpha\) ；诸 \(\delta_{\lambda}\) 是集合 \(\mathbf{N}^{(L)} = \{0\}\) 的极小元素；关系 \(\alpha \leqslant \beta\) 蕴含 \(|\alpha| \leqslant |\beta|\) ，且若 \(\alpha \leqslant \beta\) 与 \(|\alpha| = |\beta|\) 两者都成立，则 \(\alpha = \beta\) 。最后，我们记 \(\alpha! = \prod_{\lambda \in L} (\alpha(\lambda))!\) ，它有意义，因为 \(0! = 1\) 。

从第4节到第8节（含），A表示一个交换环，且除非另有说明，所考虑的代数均假定为结合的且有单位的，代数同态也假定为有单位的。

